% from main.tex
%%%%%%%%%%  scan idx 92  |  folio 79  %%%%%%%%%%
\[
  L_n\longrightarrow L_{n-1}\longrightarrow \cdots \longrightarrow L_0\longrightarrow F\longrightarrow 0,
\]
où les \(L_i\)
sont libres de type fini. Le faisceau \(\mathcal O_X\) n'étant plus supposé cohérent, disons qu'un \(\mathcal O_X\)-Module \(F\) est \emph{pseudo-cohérent} s'il vérifie la condition précédente, i.e.\ est de \(n\)-présentation finie pour tout \(n\). Alors, la notion de pseudo-cohérence possède, vis-à-vis des suites exactes courtes, la même propriété de stabilité que la notion de cohérence : si deux des termes d'une suite exacte courte de \(\mathcal O_X\)-Modules sont pseudo-cohérents, le troisième l'est aussi.

On généralise aisément la notion précédente aux complexes de faisceaux. Un \emph{complexe} \(F\) de \(\mathcal O_X\)-Modules est dit \emph{pseudo-cohérent} s'il est \(n\)-pseudo-cohérent pour tout \(n\in\mathbb Z\), « \(n\)-pseudo-cohérent » signifiant qu'il existe localement un quasi-isomorphisme \(L\to F\), où \(L\) est un complexe à degrés bornés supérieurement et à composantes libres de type fini en degré \(\geq n\). Pour un complexe \(F\) réduit au degré \(0\), il revient au même de dire que \(F\) est pseudo-cohérent en tant que complexe ou en tant que Module.

La notion de complexe pseudo-cohérent a d'excellentes vertus de stabilité, décrites dans Exp.\ I \S\S~1 et 2. D'abord c'est une notion stable par isomorphisme dans la catégorie dérivée \(\Db(X)\) de la catégorie des \(\mathcal O_X\)-Modules; qui mieux est, si deux des objets d'un triangle distingué de \(\Db(X)\) sont pseudo-cohérents, le troisième l'est aussi : en d'autres termes, la sous-catégorie pleine \(\Db(X)_{\mathrm{coh}}\) de \(\Db(X)\) formée des complexes pseudo-cohérents est une sous-catégorie triangulée. En outre, la pseudo-cohérence se conserve par image inverse et produit tensoriel dérivés. Dans le cas où \(\mathcal O_X\) est cohérent, dire qu'un complexe \(F\) est pseudo-cohérent signifie simplement que \(F\) est à cohomologie localement bornée supérieurement et que tous les \(\underline H^i(F)\) sont cohérents.

Si \(\mathcal O_X\) est un faisceau d'anneaux locaux réguliers, tout \(\mathcal O_X\)-Module
